% ============================================================================
% CHAPTER 4: IMPLEMENTATION
% ============================================================================

\chapter{IMPLEMENTATION}

\section{DEVELOPMENT ENVIRONMENT}

\subsection{Development Tools}

\begin{table}[H]
\centering
\caption{Development Tools and Versions}
\label{tab:dev_tools}
\begin{tabular}{|l|l|l|}
\hline
\textbf{Tool} & \textbf{Version} & \textbf{Purpose} \\
\hline
Node.js & 18.x+ & JavaScript runtime \\
\hline
TypeScript & 5.3.0 & Type-safe JavaScript \\
\hline
Python & 3.10+ & Backend runtime \\
\hline
FastAPI & 0.109.0 & Web framework \\
\hline
XGBoost & 2.0.3 & ML classifier \\
\hline
scikit-learn & 1.4.0 & ML utilities \\
\hline
Pydantic & 2.5.3 & Data validation \\
\hline
React & 17.0+ & UI framework (peer) \\
\hline
\end{tabular}
\end{table}

\subsection{Project Structure}

\begin{verbatim}
passiveguard/
├── frontend/                    # TypeScript SDK
│   ├── src/
│   │   ├── types/              # Type definitions
│   │   │   └── index.ts
│   │   ├── collectors/         # Data collectors
│   │   │   ├── environmental.ts
│   │   │   └── behavioral.ts
│   │   ├── utils/              # Utilities
│   │   │   ├── api.ts
│   │   │   └── fingerprint.ts
│   │   ├── core/               # Core SDK
│   │   │   └── PassiveGuard.ts
│   │   ├── hooks/              # React hooks
│   │   │   └── usePassiveGuard.ts
│   │   └── index.ts            # Entry point
│   ├── package.json
│   └── tsconfig.json
├── backend/                     # Python API
│   ├── app/
│   │   ├── api/
│   │   │   ├── models.py       # Pydantic schemas
│   │   │   └── routes.py       # API endpoints
│   │   ├── core/
│   │   │   ├── config.py       # Configuration
│   │   │   └── security.py     # JWT handling
│   │   ├── features/
│   │   │   └── extractor.py    # Feature extraction
│   │   ├── ml/
│   │   │   └── detector.py     # Bot detector
│   │   └── main.py             # FastAPI app
│   ├── ml_training/
│   │   └── train_model.py      # Model training
│   ├── models/                  # Trained models
│   │   └── bot_detector.joblib
│   └── requirements.txt
└── demo/
    └── index.html              # Demo application
\end{verbatim}

\section{FRONTEND SDK IMPLEMENTATION}

\subsection{Type Definitions}

The SDK defines comprehensive TypeScript interfaces for type safety:

\begin{verbatim}
// Core configuration interface
interface PassiveGuardConfig {
  apiEndpoint: string;
  siteKey: string;
  debug?: boolean;
  collectionTimeout?: number;
  enableBehavioral?: boolean;
  enableEnvironmental?: boolean;
  privacyMode?: boolean;
  onChallenge?: (challenge: Challenge) => Promise<ChallengeResponse>;
  onVerify?: (result: VerificationResult) => void;
}

// Verification result interface
interface VerificationResult {
  isHuman: boolean;
  confidence: number;
  riskScore: number;
  token: string;
  timestamp: number;
  challengeRequired: boolean;
  challenge?: Challenge;
  metadata: VerificationMetadata;
}
\end{verbatim}

\subsection{Environmental Collector}

The EnvironmentalCollector class gathers browser and device characteristics:

\begin{verbatim}
class EnvironmentalCollector {
  async collect(): Promise<EnvironmentalData> {
    const [canvasFp, webglFp, audioFp] = await Promise.all([
      this.getCanvasFingerprint(),
      this.getWebGLFingerprint(),
      this.getAudioFingerprint()
    ]);
    
    return {
      browser: this.getBrowserInfo(),
      screen: this.getScreenInfo(),
      hardware: this.getHardwareInfo(),
      features: this.getFeatureDetection(),
      fingerprints: {
        canvas: this.hashFingerprint(canvasFp),
        webgl: this.hashFingerprint(webglFp),
        audio: this.hashFingerprint(audioFp)
      },
      timezone: this.getTimezoneInfo()
    };
  }
}
\end{verbatim}

\subsection{Canvas Fingerprinting Implementation}

\begin{verbatim}
private getCanvasFingerprint(): string {
  const canvas = document.createElement('canvas');
  const ctx = canvas.getContext('2d');
  if (!ctx) return '';

  canvas.width = 200;
  canvas.height = 50;

  // Draw text with specific styling
  ctx.textBaseline = 'top';
  ctx.font = '14px Arial';
  ctx.fillStyle = '#f60';
  ctx.fillRect(125, 1, 62, 20);

  ctx.fillStyle = '#069';
  ctx.fillText('PassiveGuard', 2, 15);
  
  ctx.fillStyle = 'rgba(102, 204, 0, 0.7)';
  ctx.fillText('Bot Detection', 4, 17);

  return canvas.toDataURL();
}
\end{verbatim}

\subsection{Behavioral Collector}

The BehavioralCollector tracks user interactions:

\begin{verbatim}
class BehavioralCollector {
  private mouseMovements: MouseEvent[] = [];
  private keyPresses: KeyboardEvent[] = [];
  private scrollEvents: ScrollEvent[] = [];
  
  start(): void {
    document.addEventListener('mousemove', this.handleMouseMove);
    document.addEventListener('keydown', this.handleKeyDown);
    document.addEventListener('keyup', this.handleKeyUp);
    document.addEventListener('scroll', this.handleScroll);
  }
  
  private handleMouseMove = (e: MouseEvent): void => {
    this.mouseMovements.push({
      x: e.clientX,
      y: e.clientY,
      timestamp: Date.now()
    });
    // Keep last 1000 movements
    if (this.mouseMovements.length > 1000) {
      this.mouseMovements.shift();
    }
  };
  
  getData(): BehavioralData {
    return {
      mouse: this.calculateMouseStats(),
      keyboard: this.calculateKeyboardStats(),
      scroll: this.calculateScrollStats(),
      timing: this.calculateTimingStats()
    };
  }
}
\end{verbatim}

\subsection{Mouse Statistics Calculation}

\begin{verbatim}
private calculateMouseStats(): MouseStats {
  const movements = this.mouseMovements;
  if (movements.length < 2) return this.getEmptyMouseStats();

  const velocities: number[] = [];
  let totalDistance = 0;
  let directionChanges = 0;
  let pauseCount = 0;

  for (let i = 1; i < movements.length; i++) {
    const dx = movements[i].x - movements[i-1].x;
    const dy = movements[i].y - movements[i-1].y;
    const dt = movements[i].timestamp - movements[i-1].timestamp;
    
    const distance = Math.sqrt(dx*dx + dy*dy);
    totalDistance += distance;
    
    if (dt > 0) {
      velocities.push(distance / dt);
    }
    
    if (dt > 100) pauseCount++;  // Pause > 100ms
  }

  return {
    totalDistance,
    averageVelocity: mean(velocities),
    velocityStd: std(velocities),
    maxVelocity: Math.max(...velocities),
    directionChanges,
    pauseCount
  };
}
\end{verbatim}

\subsection{PassiveGuard Core Class}

The main SDK class orchestrates collection and verification:

\begin{verbatim}
class PassiveGuard {
  private config: PassiveGuardConfig;
  private environmentalCollector: EnvironmentalCollector;
  private behavioralCollector: BehavioralCollector;
  private apiClient: ApiClient;
  
  async init(): Promise<void> {
    if (this.config.enableBehavioral) {
      this.behavioralCollector.start();
    }
    if (this.config.enableEnvironmental) {
      this.environmentalData = await 
        this.environmentalCollector.collect();
    }
    this.initialized = true;
  }
  
  async verify(): Promise<VerificationResult> {
    const request: VerificationRequest = {
      siteKey: this.config.siteKey,
      environmental: this.environmentalData,
      behavioral: this.behavioralCollector.getData(),
      requestId: this.generateRequestId(),
      timestamp: Date.now()
    };
    
    return this.apiClient.verify(request);
  }
}
\end{verbatim}

\subsection{React Hook Implementation}

\begin{verbatim}
function usePassiveGuard(options: UsePassiveGuardOptions) {
  const [isInitialized, setIsInitialized] = useState(false);
  const [isVerifying, setIsVerifying] = useState(false);
  const [result, setResult] = useState<VerificationResult | null>(null);
  const [error, setError] = useState<Error | null>(null);
  
  const guardRef = useRef<PassiveGuard | null>(null);
  
  useEffect(() => {
    guardRef.current = new PassiveGuard({
      apiEndpoint: options.apiEndpoint,
      siteKey: options.siteKey,
      ...options
    });
    
    guardRef.current.init()
      .then(() => setIsInitialized(true))
      .catch(setError);
      
    return () => guardRef.current?.destroy();
  }, [options.apiEndpoint, options.siteKey]);
  
  const verify = useCallback(async () => {
    setIsVerifying(true);
    try {
      const result = await guardRef.current!.verify();
      setResult(result);
      return result;
    } finally {
      setIsVerifying(false);
    }
  }, []);
  
  return { isInitialized, isVerifying, result, error, verify };
}
\end{verbatim}

\section{BACKEND API IMPLEMENTATION}

\subsection{FastAPI Application}

\begin{verbatim}
from fastapi import FastAPI
from fastapi.middleware.cors import CORSMiddleware

app = FastAPI(
    title="PassiveGuard API",
    description="ML-based passive bot detection",
    version="1.0.0"
)

app.add_middleware(
    CORSMiddleware,
    allow_origins=["*"],
    allow_credentials=True,
    allow_methods=["*"],
    allow_headers=["*"]
)

@app.on_event("startup")
async def startup():
    # Load ML model
    global detector
    detector = BotDetector()
\end{verbatim}

\subsection{Pydantic Models}

\begin{verbatim}
class EnvironmentalData(BaseModel):
    browser: BrowserInfo
    screen: ScreenInfo
    hardware: HardwareInfo
    features: FeatureDetection
    fingerprints: Fingerprints
    timezone: TimezoneInfo

class BehavioralData(BaseModel):
    mouse: MouseData
    keyboard: KeyboardData
    scroll: ScrollData
    timing: TimingData

class VerificationRequest(BaseModel):
    siteKey: str
    environmental: EnvironmentalData
    behavioral: BehavioralData
    requestId: str
    timestamp: int

class VerificationResponse(BaseModel):
    isHuman: bool
    confidence: float
    riskScore: float
    token: str
    timestamp: int
    challengeRequired: bool
\end{verbatim}

\subsection{Verification Endpoint}

\begin{verbatim}
@router.post("/verify", response_model=VerificationResponse)
async def verify(request: VerificationRequest):
    # Extract features
    features = extractor.extract(
        request.environmental,
        request.behavioral
    )
    
    # Get prediction
    prediction = detector.predict(features)
    
    # Generate token
    token = create_verification_token(
        is_human=prediction.is_human,
        confidence=prediction.confidence,
        request_id=request.requestId
    )
    
    return VerificationResponse(
        isHuman=prediction.is_human,
        confidence=prediction.confidence,
        riskScore=prediction.risk_score,
        token=token,
        timestamp=int(time.time() * 1000),
        challengeRequired=prediction.needs_challenge
    )
\end{verbatim}

\subsection{Feature Extractor}

\begin{verbatim}
class FeatureExtractor:
    def extract(self, env: EnvironmentalData, 
                beh: BehavioralData) -> np.ndarray:
        features = []
        
        # Environmental features
        features.extend([
            env.screen.width / 1920,  # Normalized
            env.screen.height / 1080,
            env.hardware.deviceMemory or 4,
            env.hardware.hardwareConcurrency or 4,
            1 if env.features.webgl else 0,
            1 if env.browser.webdriver else 0,
            # ... more features
        ])
        
        # Behavioral features
        mouse = beh.mouse.movementStats
        features.extend([
            mouse.totalDistance,
            mouse.averageVelocity,
            mouse.velocityStd,
            mouse.directionChanges,
            mouse.pauseCount,
            # ... more features
        ])
        
        return np.array(features).reshape(1, -1)
\end{verbatim}

\section{ML MODEL TRAINING}

\subsection{Synthetic Data Generation}

Training data is generated by simulating human and bot behaviors:

\begin{verbatim}
def generate_human_sample():
    return {
        'mouse_velocity_mean': np.random.normal(150, 50),
        'mouse_velocity_std': np.random.normal(80, 30),
        'mouse_direction_changes': np.random.randint(10, 100),
        'mouse_pause_count': np.random.randint(5, 30),
        'keyboard_dwell_mean': np.random.normal(100, 30),
        'scroll_smoothness': np.random.uniform(0.6, 1.0),
        'timing_first_interaction': np.random.normal(2000, 1000),
        'env_webdriver': 0,
        # ... more features
    }

def generate_bot_sample():
    return {
        'mouse_velocity_mean': np.random.uniform(0, 50),
        'mouse_velocity_std': np.random.uniform(0, 10),
        'mouse_direction_changes': np.random.randint(0, 5),
        'mouse_pause_count': np.random.randint(0, 2),
        'keyboard_dwell_mean': np.random.uniform(50, 80),
        'scroll_smoothness': np.random.uniform(0.0, 0.3),
        'timing_first_interaction': np.random.uniform(0, 500),
        'env_webdriver': np.random.choice([0, 1], p=[0.3, 0.7]),
        # ... more features
    }
\end{verbatim}

\subsection{Model Training}

\begin{verbatim}
def train_model():
    # Generate training data
    human_samples = [generate_human_sample() for _ in range(10000)]
    bot_samples = [generate_bot_sample() for _ in range(10000)]
    
    X = pd.DataFrame(human_samples + bot_samples)
    y = np.array([1] * 10000 + [0] * 10000)  # 1=human, 0=bot
    
    # Train-test split
    X_train, X_test, y_train, y_test = train_test_split(
        X, y, test_size=0.2, random_state=42
    )
    
    # Train XGBoost
    model = XGBClassifier(
        n_estimators=100,
        max_depth=6,
        learning_rate=0.1,
        subsample=0.8,
        colsample_bytree=0.8,
        eval_metric='logloss',
        early_stopping_rounds=10
    )
    
    model.fit(
        X_train, y_train,
        eval_set=[(X_test, y_test)],
        verbose=True
    )
    
    # Save model
    joblib.dump(model, 'models/bot_detector.joblib')
\end{verbatim}

\section{SECURITY IMPLEMENTATION}

\subsection{JWT Token Generation}

\begin{verbatim}
def create_verification_token(is_human: bool, 
                              confidence: float,
                              request_id: str) -> str:
    payload = {
        "sub": request_id,
        "is_human": is_human,
        "confidence": confidence,
        "iat": datetime.utcnow(),
        "exp": datetime.utcnow() + timedelta(minutes=5)
    }
    return jwt.encode(payload, SECRET_KEY, algorithm="HS256")

def validate_token(token: str) -> dict:
    try:
        payload = jwt.decode(token, SECRET_KEY, algorithms=["HS256"])
        return {"valid": True, **payload}
    except jwt.ExpiredSignatureError:
        return {"valid": False, "error": "Token expired"}
    except jwt.InvalidTokenError:
        return {"valid": False, "error": "Invalid token"}
\end{verbatim}

\subsection{Fingerprint Hashing}

\begin{verbatim}
// Frontend - TypeScript
private hashFingerprint(data: string): string {
  // Using SubtleCrypto for SHA-256
  const encoder = new TextEncoder();
  const dataBuffer = encoder.encode(data);
  
  return crypto.subtle.digest('SHA-256', dataBuffer)
    .then(hash => {
      return Array.from(new Uint8Array(hash))
        .map(b => b.toString(16).padStart(2, '0'))
        .join('');
    });
}
\end{verbatim}

\section{DEMO APPLICATION}

The demo application provides a working example of SDK integration:

\begin{verbatim}
<!DOCTYPE html>
<html>
<head>
  <title>PassiveGuard Demo</title>
</head>
<body>
  <div id="app">
    <h1>PassiveGuard Bot Detection Demo</h1>
    <div id="status">Initializing...</div>
    <button id="verify-btn" disabled>Verify Human</button>
    <div id="result"></div>
  </div>
  
  <script type="module">
    import { PassiveGuard } from './passiveguard-sdk.js';
    
    const guard = new PassiveGuard({
      apiEndpoint: 'http://localhost:8000',
      siteKey: 'demo_site',
      debug: true
    });
    
    await guard.init();
    document.getElementById('status').textContent = 'Ready';
    document.getElementById('verify-btn').disabled = false;
    
    document.getElementById('verify-btn').onclick = async () => {
      const result = await guard.verify();
      document.getElementById('result').innerHTML = `
        <p>Human: ${result.isHuman}</p>
        <p>Confidence: ${(result.confidence * 100).toFixed(1)}%</p>
      `;
    };
  </script>
</body>
</html>
\end{verbatim}

